\scanpage{23}
% archive sheet number: 22

\unsure{6|§|(} Automorphismes de $X$ $(=\mathrm P(\unsure{E|\mathcal E}))$ induisant l'identité sur $C$.

Ils correspondent aux pseudo-reflexions de $E$ qui sont l'identité sur $V$, donc de la forme
\[
u=\mathrm{id}+\pi\otimes a\tag*{(1)}
\]
$a$ section de $E$. Si on écrit que $u$ commute aux $\tau_i$, cela s'écrit $\langle a,a'_i\rangle=0$ $(0\leq i\leq n)$ i.e. $a$ de la forme $\lambda\widetilde o$
\[
u=\mathrm{id}+\pi\otimes\lambda\widetilde o\tag*{(2)}
\]

Écrivons que $u$ est orthogonale, i.e.
\[
\underbrace{f(\deleted{\unsure{a|\lambda}}\lambda\widetilde o)}_{-\lambda^2\sigma'\rho'}\pi+
\underbrace{\psi(\lambda\widetilde o)}_{-\lambda\sigma\pi}=0
\]

ou encore
\[
\lambda(\lambda\sigma'\rho'+\sigma)=0\tag*{(3)}
\]
ce qui suivant la parité de $n$ s'écrit
\[
\left\{\begin{aligned}
 &n\text{ pair }(\rho=2\rho',\underline{\sigma=\sigma'})&
 \lambda\sigma'(\lambda\rho'+1)&=0\\
 &n\text{ impair }(\underline{\rho=\rho'},\sigma=2\sigma')&
 \lambda\sigma'(\lambda\rho^{\deleted{\unsure{\prime}}}+2)&=0
\end{aligned}\right.\tag*{(4)}
\]

D'autre part, la \deleted{fait} \inserted{condition} que $u$ soit une \uline{reflexion} s'écrit $\langle\lambda\widetilde o,\pi\rangle=-2$ i.e.
\[
\lambda\rho=-2\tag*{(5)}
\]

\deleted{\unsure{$n$ pair}} S

\uline{Proposition} 1) Si $\sigma'=0$, \unsure{alors|tous|toutes} les \deleted{\unsure{auto|autres}} \deleted{\unsure{qui}} pseudo-reflexions \deleted{telles que} \inserted{commutant à $G$ et} fixes sur $V$ sont orth.

2) Si $\sigma'$ inv, \deleted{une}\inserted{la} ps. reflexion $u=\mathrm{id}+\lambda\pi\otimes\widetilde o$
